\documentclass[10pt,letterpaper]{article}
\usepackage[margin=1in,textwidth=6.4in]{geometry}
\usepackage{amsmath,amssymb,graphicx,booktabs,longtable,xcolor}
\usepackage[breaklinks=true,hidelinks]{hyperref}
\newcommand{\placeholder}[1]{\textcolor{orange}{#1}}
\newcommand{\PaperSourceRevision}{2bd3103a90b5}
\newcommand{\MaxSpeckleNoise}{$3.5\times10^{-4}$}
\newcommand{\MaxTotalNoise}{$7.5\times10^{-3}$}
\newcommand{\MaxEmissivityTotalNoise}{$0.15$}
\newcommand{\ReferenceSpeckleNoise}{$9.0\times10^{-8}$}
\newcommand{\PhysicalCaseCount}{$520$}
\newcommand{\PhysicalRms}{$5$}
\newcommand{\EmittingPathPc}{$5$}
\newcommand{\UnitEmissivity}{$1$}
\newcommand{\MeanCaseCount}{$13$}
\newcommand{\MaxMeanDepth}{$0.66$}
\newcommand{\LunarFaradayBound}{$2.6\times10^{-6}$}
\newcommand{\LunarPhaseBound}{$4.7\times10^{-3}$}
\newcommand{\ParentResidualDof}{$1$}
\newcommand{\ZoomResidualDof}{$190$}
\newcommand{\ResidualCaseCount}{$40$}

\title{Conditional Faraday polarization predictions for LuSEE-Night}
\author{Implementation and numerical validation report}
\date{Source revision \texttt{\PaperSourceRevision}}
\begin{document}
\maketitle

\noindent\fbox{\parbox{0.95\linewidth}{\textbf{Status.} Conditional model
calculation and simulated observations. The comparison target is future
LuSEE-Night data. Local polarized emissivity, small-scale Faraday statistics,
receiver current noise and fine-channel beam errors remain assumptions.
This report makes no calibrated detection or launch-date forecast.}}

\section{Question and measured predictions}

The calculation predicts how much polarization survives a distribution of
emitting depths, which part passes the loaded four-port instrument, and
which part remains after channel averaging and smooth subtraction.
Its outputs are means and joint covariances of calibrated Stokes spectra,
parent/zoom residual variance, frequency correlations, and differences
between repeated observations at the same lunar sidereal phase. Simulated
observations exercise the same calibrated-data comparison interface that
can later accept mission spectra with their actual calibration and responses.

Across the full signed-column grid at \PhysicalRms\ K intrinsic RMS,
the largest residual speckle to radiometer RMS ratio is \MaxSpeckleNoise;
including smooth polarization gives \MaxTotalNoise. These amplitude ratios
are not detection significances. The exported decision records also retain
signal transmission, residual leakage, beam ablation mismatch, interpolation
spread and numerical error. The interpretation is conditional on every
input in Table~\ref{tab:inputs}.

\section{Physical model and local inputs}

Write the COSMO polarization as
\begin{equation}
 P(\hat n,x)=P_0(\hat n)\int_0^1 w(f)
       \exp[2ifx\Phi(\hat n)]\,df,\qquad x=\lambda^2.
\end{equation}
Here $f$ is the accumulated rotating-column fraction: $f=0$ is unrotated
and $f=1$ recovers the foreground-screen pilot. The prescribed fractional
screens, uniform slab and supported Gaussian all share one column along
each sightline. They bracket depth mixing with one intrinsic polarization
angle; they do not represent independent emitting layers or a complete
three-dimensional turbulent medium.

For $A=\int w(f)e^{2ifx\Phi}df$, use $A=m+\delta A$ with
$m=E_\Phi A$. The smooth field is $mP_0$ and the speckle is
$\delta A P_0$. An isotropic intrinsic sky can have zero ensemble mean
and still have smooth power. Here ``smooth'' means the conditional-mean
field; at a large signed mean column it can rotate rapidly with frequency.
Both the conjugated covariance and the
nonconjugated relation are retained before mapping to real readouts.
The single-depth raw relation at a coincident point is
$e^{4i\mu x-8\sigma_\Phi^2x^2}$; this expression is not imposed on a slab
or an angularly integrated centered covariance.

\begin{table}[htbp]\centering\small
\caption{Adopted inputs and their status.}\label{tab:inputs}
\begin{tabular}{p{0.22\linewidth}p{0.26\linewidth}p{0.43\linewidth}}\toprule
Quantity & Input & Source/status\\\midrule
Reference / volume RMS & 1 / 5 K at 30 MHz & Illustrative reference / emissivity times full path\\
Volume emissivity / path & 1 K pc$^{-1}$ / 5 pc & Assumed polarized RMS; depth integration supplies suppression\\
Intrinsic angular shape & ell scale 8; EE 0.8, EB 0.1 & Gaussian-decaying spectrum; BB is the remaining fraction\\
Intrinsic frequency shape & Index -2.8; turnover 15 MHz & Assumed absorbed power law, not a measured local spectrum\\
Angular RM structure & theta0=0.01 rad; alpha=1.667 & Assumed powered-chord covariance, fixed across mean columns\\
Local dispersion & 0.002 / 0.02 rad m$^{-2}$ & Conditional angular screens; no column-map substitution\\
Voltage scales & 2.2 / 3 / 5 nV Hz$^{-1/2}$ & Adopted EM/requirement scales; \href{https://doi.org/10.1029/2023RS007925}{Tamura et al.}\\
Current scales & 0.03 / 0.1 / 0.3 pA Hz$^{-1/2}$ & Assumed; independent input sources\\
Moon / loss temperatures & 100 K / 100 K & Illustrative nightside values\\
Emissivity audit & Total I, paths below 1 kpc & \href{https://arxiv.org/html/1811.10310v1}{Polderman et al.}; no local polarized prior\\
Nearby RM/DM & Historical J0437-4715 values & \href{https://arxiv.org/pdf/1105.4213}{Yan et al.}; path averages\\
Lunar plasma & Assumed density/field/path envelope & \href{https://doi.org/10.1029/2011JA017141}{Ando et al.}; no universal nightside limit\\
Observing duty & 50 percent & Assumed as in \href{https://arxiv.org/html/2508.16773v2}{Camacho et al.}; actual-site visibility calculated\\
\bottomrule\end{tabular}\end{table}


The signed-column calculation uses \MeanCaseCount\ mean columns between
minus \MaxMeanDepth\ and plus \MaxMeanDepth rad m$^{-2}$. They follow from
extending each adopted density/field pair uniformly over \EmittingPathPc\ pc
with both field signs and the stated projection fractions. This is a
sensitivity grid, not a measured distribution of local columns. Its intrinsic
normalization is \UnitEmissivity\ K pc$^{-1}$ times that full path, giving
\PhysicalRms\ K RMS at the reference frequency. The depth integral supplies
depolarization; the emitting volume is not truncated again at a phase horizon.
The measured-spectra illustration separately retains the original one-kelvin
reference case and its preselected mean column.

The horizon table uses a uniform signed rate,
$d\Phi/ds=0.812 n_e B_\parallel$ in rad m$^{-2}$ pc$^{-1}$, with the
electron density in cm$^{-3}$ and field in microgauss. A one-radian
complex-polarization tolerance is $2|\Phi|\lambda^2=1$; a two-radian
tolerance gives twice the listed distance. Zero line-of-sight field gives
an unbounded horizon. These arithmetic scenarios do not determine a local
$\sigma_\Phi$ or angular structure function.

\begin{table}[htbp]
\centering\small
\caption{Conditional aligned-field one-radian horizons (pc). Zero field projection is unbounded in every scenario.}
\label{tab:horizon}
\begin{tabular}{llll}
\toprule
Scenario & 10 MHz & 30 MHz & 50 MHz\\
\midrule
LIC\_model26 & $4.2\times10^{-3}$ & $0.038$ & $0.1$\\
LHB\_density\_low\_assumed\_field & $0.2$ & $1.8$ & $4.9$\\
LHB\_density\_high\_assumed\_field & $0.15$ & $1.4$ & $3.8$\\
\bottomrule
\end{tabular}
\end{table}


The Local Interstellar Cloud model pair is taken from Tables 2 and 5 of
\href{https://astro.uchicago.edu/~frisch/2008SlavinFrisch.pdf}{Slavin and
Frisch (2008)}. The Local Hot Bubble density range follows the adopted values
in \href{https://arxiv.org/html/2410.23345v1}{Yeung et al. (2024)};
its field strength here is assumed. These density ranges do not determine
the line-of-sight field orientation, so the phase horizon cannot establish
a universal local-emission volume.

The nearby-pulsar DM/RM constraint and HII-region synchrotron emissivity
average over paths much longer than the phase-limited emitting volume.
Their source and selection audit is recorded in
\nolinkurl{docs/literature/local-faraday-horizon.md}. They cannot identify
local polarized variance. The lunar plasma calculation likewise uses an
explicit aligned-field density/path envelope: its largest adopted Faraday
bound is \LunarFaradayBound\ rad m$^{-2}$, giving complex-polarization
phase \LunarPhaseBound\ rad at the lowest modeled frequency. The
occultation literature does not establish this as a universal nightside
upper limit. Emissivity, emitting-depth distribution, angular coherence and
Faraday variance remain degenerate in a single measured floor.

\section{Instrument, channels and common estimator}

All production polarization predictions use harmonic kernels from the
as-built four-port response, with the beam and receiver loading frozen at
each native band center. Signal-only calculations have zero Moon and metal
temperatures. The calibrated complex X/Y combinations and one-kelvin loaded
enclosure normalization are shared by every component.
The I-only leakage controls instead evaluate the actual native-frequency
beam and loading for the as-built, C4-symmetrized and diagonal-ZA responses.
Each uses the as-built band's calibration, so their differences include
the mismatch relevant to one model correcting another.

Haslam synchrotron and CMB intensity maps are transformed at their native
resolution, with no map degradation. Native response samples define linear
and PCHIP fine-band interpolants. Differences between those interpolants
are a model spread; they do not measure beam structure inside a native
frequency interval. Fitting a frozen kernel would give an uninformative
leakage-smoothness test.

The actual power-response matrix $W$ carries means and both moment types
through folded parent/zoom responses. A constant and slope in physical fine
frequency pass through $W$ before the fit. The resulting residual map $R$
acts on leakage, smooth polarization, speckle and noise identically. For
structured phase coefficients, applying $R$ before forming covariance avoids
subtraction of nearly equal large terms. Three parents retain
\ParentResidualDof\ residual degree of freedom; the zoom fit retains
\ZoomResidualDof. The saturated parent quadratic is rejected.

Q/U subtraction is performed in kelvin before division by a common
channelized intensity reference. The exported estimated-I alternative
retains denominator variance and its correlations with Q/U through the
full Jacobian. Low intensity signal-to-noise is rejected as outside that
linearization. A zero comparison denominator remains undefined.

\begin{table}[htbp]
\centering\small
\caption{Noise-equivalent bandwidth from actual folded power weights. Ranges span channel positions, not uncertainty intervals.}
\label{tab:channels}
\begin{tabular}{llll}
\toprule
Type & Bins & Nominal spacing (Hz) & Effective bandwidth (Hz)\\
\midrule
Parent & 3 & 25000 & $2.473\times10^{4}$\\
Zoom & 192 & 390.625 & $[581.4,1.206\times10^{3}]$\\
\bottomrule
\end{tabular}
\end{table}

\begin{table}[htbp]
\centering\small
\caption{Nominal narrow-channel top-hat first-zero Faraday depth $\pi/\Delta(\lambda^2)$, in rad m$^{-2}$. This is a resolution scale; actual predictions use the measured folded power weights rather than this top-hat approximation.}
\label{tab:resolution}
\begin{tabular}{lll}
\toprule
MHz & Parent & Zoom\\
\midrule
10 & $0.7$ & $45$\\
30 & $19$ & $1.2\times10^{3}$\\
50 & $87$ & $5.6\times10^{3}$\\
\bottomrule
\end{tabular}
\end{table}


Compact all-product exports retain every real science channel in the pinned
16-product order. Reprojection to calibrated Stokes checks each Hermite
mode, I-only means, and raw/residual radiometer covariances with one common
normalization. The complete channel covariance remains factorized until
requested; selected parent/zoom blocks are also exported explicitly.

\section{Receiver noise and observing exposure}

The total electric-field spectral covariance sums loaded sky, Moon,
antenna-loss and amplifier terms once. Amplifier voltage and current
sources propagate through the full coupled antenna/load network; impedances
alone do not supply amplifier noise. The public engineering-model voltage
measurements set adopted reference scales, while current noise, input-source
correlations and thermal temperatures are explicitly assumed. Missing
legacy SPICE files are not replaced by zero noise.

For proper complex voltages, real quadratic readouts obey
\begin{equation}
 \mathrm{Cov}(q_p,q_q)=\mathrm{Tr}(A_p S A_q S)/(B\tau).
\end{equation}
Fine-frequency weighting propagates cross-product and cross-channel
covariance, including overlapping and folded responses and frequency-varying
system power. This uses the stated independent fine-spectrum and
long-integration approximation; a detailed temporal filter model would be
an additional input. Receiver bias spectra are exported; the sky-I examples
assume their mean has been calibrated. Polarization is small compared with
I in the normalized illustrations, allowing the noise calculation to use
I-only sky power.

Every phase-bin exposure counts complete bins that are nightside at the
actual lunar site, at both edges and at the center. A declared solar-altitude
margin and electronics duty apply before averaging. MEPA harmonic kernels
receive the corresponding exact boxcar sinc factor on both signal and
leakage paths. The calendar origin is a simulation coordinate, not a launch
date. An unavailable phase has no reported noise prediction.

\begin{table}[htbp]
\centering\small
\caption{Illustrative actual-site nightside exposure; the electronics duty is applied to eligible complete bins.}
\label{tab:exposure}
\begin{tabular}{llll}
\toprule
Lunations & Sidereal turn & Visits & Integration (s)\\
\midrule
1 & 0.0 & 1 & $1.153\times10^{3}$\\
1 & 0.25 & 0 & Unavailable\\
24 & 0.0 & 14 & $1.614\times10^{4}$\\
24 & 0.25 & 14 & $1.614\times10^{4}$\\
\bottomrule
\end{tabular}
\end{table}


\section{Numerical checks and scope}

The implementation retains all helicities, pair products and the complex
relation. Depth-averaged Hermite coefficients use a bounded omitted-mode
tail and independently refined depth quadrature. Spin/angular, beam and
frequency refinements are checked separately for the unfiltered mixed-depth
cases. The residual production stage contains \ResidualCaseCount\ cases;
their common-fit covariance positivity and subtraction identity are checked
separately. The expanded signed-column grid adds \PhysicalCaseCount\ cases
with exact diagonal contractions of the same complete complex operator.
Frequency/depth refinement bounds the full Q/U covariance change; angular
and band-limit refinements cover every reported marginal. Full joint checks
also cover eight predeclared extreme-column screen/slab cases, and saved
coefficients allow a complete covariance export for any grid case. Tiny
residual power must be compared with its exported numerical error before
interpretation.

The deterministic narrow-window spherical control recovers the earlier
flat-plane PR~7/8 source geometry. The isotropic random intrinsic E/B sky
used in production remains a different ensemble. An independent
bivariate-Gaussian channel control checks the global variance mixture;
matching its RM covariance with a Gaussian screen does not match its phase
statistics. Strong-screen controls use the direct frequency-pair solver
when a Hermite expansion would require excessive order, and test the scalar
floor against the full loaded readout normalization.

Independent voltage ensembles and exact Gaussian fourth-moment cubature
check radiometer normalization and sampling uncertainty. Deliberate faults
test the covariance legs, phases, folds, depth mixing, intensity ratios,
noise loading and repeat scaling. Repeated visits reuse one sky and screen;
only independent measurement noise decreases. The Gaussian visibility
draws illustrate conditional second moments. They do not establish
Gaussianity of every weak-rotation or few-cell sky. Separate band
illustrations use independent seeds and do not supply a joint cross-band
sky realization; within each band, channels, phases and repeats share the
same draw.

\section{Figures and comparison with data}

\begin{figure}[p]\centering
\includegraphics[width=6.4in]{figures/faraday_instrument_leakage.pdf}
\caption{Instrument and diffuse leakage. I-only native-frequency Haslam+CMB harmonic predictions at sidereal phase 0 through the as-built, C4-symmetrized and diagonal-ZA responses. Every variant uses the as-built band's fixed zenith calibration and enclosure normalization, plus identical MEPA time averaging. Curves show actual native frequencies; fine-channel linear/PCHIP leakage and their common-fit residuals are model interpolants exported in the result record. Variant differences are a model spread, not a beam-error probability distribution. These diffuse-sky leakages need not vanish after zenith calibration.}
\end{figure}
\begin{figure}[p]\centering
\includegraphics[width=6.4in]{figures/paper_depth_depolarization.pdf}
\caption{Depth-profile benchmark. Conditional flat-plane depth-profile moments. Smooth \textbar{}E[P]\textbar{} squared and centered speckle power are divided by the same unrotated beam-weighted reference. The Gaussian profile is supported on f in [0,1]; fractions 0/0.6/1 retain the original screen limits. Monochromatic points use 10/30/50 MHz; lines guide the eye. Triangles mark values below the 1e-12 display floor, including exact zeros and underflow. The source/beam geometry and local screen parameters remain conditional. This independent benchmark includes no real four-port beam, channel averaging, leakage or receiver noise; those enter the subsequent instrument predictions.}
\end{figure}
\begin{figure}[p]\centering
\includegraphics[width=6.4in]{figures/faraday_redistribution_spin2.pdf}
\caption{Real-instrument redistribution. Conditional as-built loaded four-port predictions at the reference MEPA sidereal phase, using the two illustrative local Gaussian columns and the five shared-column emitting-depth profiles. The first two rows show Var(Q)+Var(U) for the phase-mean and centered terms, divided by the same unrotated-beam reference power. The bottom row is sqrt[Var(Q)+Var(U)] in K. Points use their own native 10/30/50 MHz beams; connecting segments guide the eye. Triangles mark zero/underflow values at the 1e-16 display floor. The intrinsic spin-2 RMS is 1 K at 30 MHz with assumed E/B spectra, beta=-2.8 and a 15 MHz turnover. This is an illustrative brightness law, especially below 20 MHz. Slab/Gaussian emission depths share one column rather than independent rotating layers. These are model second moments, without I-only leakage or measurement noise.}
\end{figure}
\begin{figure}[p]\centering
\includegraphics[width=6.4in]{figures/faraday_measured_spectra.pdf}
\caption{Simulated measured spectra. Simulated observations, not mission data: as-built calibrated Q/I and U/I at preselected sidereal phase 0, uniform emitting slab and sigma\_phi=0.02 rad/m$^2$. Zoom curves retain all folded images; open circles are three parent measurements. A constant and physical-frequency slope pass through each channel response before fitting. Residuals are subtracted in K, then divided by the same fixed channelized sky-I reference as the components. Estimated-I covariances, including I-Q/U correlations, are exported. The paired I-only control shares the noise draw. The last row shows the underlying smooth and speckle polarization draws; only noise changes between repeated visits. Intrinsic RMS is 1 K at 30 MHz; local structure, turnover, current noise and temperatures are assumed. Voltage/current references are 3 nV/sqrt(Hz), 0.1 pA/sqrt(Hz). Exposure spans 24 illustrative lunations at 50\% duty. The colored band shows marginal conditional sky-plus-noise RMS around the residual I-only model (dashed). Native interpolation limits fine-channel leakage interpretation.}
\end{figure}
\begin{figure}[p]\centering
\includegraphics[width=6.4in]{figures/faraday_channel_covariance.pdf}
\caption{Frequency correlations. Conditional slab-signal Q correlations before and after the common smooth fit, on all 192 nominal zoom frequencies, for the same model and time bin as the measured-spectra figure. Before-fit images show 1-r on a common log scale with floor 1e-12; after-fit images show signed r. The last row retains individual Q and U correlations to the named center-frequency bin; unlike frequency pairs are not averaged into a stationary estimate. The noise-control line includes channel folds, overlap and the same fit. Full Q-U cross covariance and both parent/zoom blocks are exported. Repeated-cycle differences and their radiometer prediction are shown in faraday\_repeatability.pdf. These conditional correlations are not a Faraday identification or an observed covariance estimate.}
\end{figure}
\begin{figure}[p]\centering
\includegraphics[width=6.4in]{figures/faraday_channel_covariance_amplitude.pdf}
\caption{Q/U covariance amplitudes. Conditional total polarized-signal Q and U covariances on all 192 nominal zoom frequencies, before and after the common fit, for the same reference slab and phase as the measured spectra. The upper four panels share one positive logarithmic scale; the lower four share one signed linear scale. Each is in K squared at the common loaded calibration plane. The before/after groups use different scales because the fit removes most of this coherent reference signal. These are ensemble signal covariances, not covariances estimated from a single spectrum. The separate noise control and full Q-U, parent/zoom and complex C/J blocks are exported; normalized correlations are shown in the companion.}
\end{figure}
\begin{figure}[p]\centering
\includegraphics[width=6.4in]{figures/faraday_separability.pdf}
\caption{Conditional separability. Conditional residual Q/U RMS divided by the same fitted radiometer RMS over all signed uniform mean columns from the adopted density/field/path pairs. Every plotted mean includes both assumed dispersions, Gaussian and global-variance-mixture statistics, all five depth profiles and nine receiver-noise combinations. Solid points are maximum speckle ratios; shading spans positive speckle ratios. The unrotated f=0 case has exactly zero speckle. Dashed open squares give maximum total (smooth plus speckle) ratios, including f=0. These are conditional model ranges, not confidence intervals. The horizontal dotted line is equal RMS, not a detection threshold. The full arrays retain both phases and 1/24-lunation availability; this figure uses phase 0 and 24 illustrative lunations. Intrinsic RMS is 5 K at 30 MHz, the assumed volume-averaged polarized emissivity times the full 5 pc path. The depth integral supplies depolarization; no second horizon cut is made. Alternative emissivity normalizations are tabulated separately and scale amplitudes linearly.}
\end{figure}
\begin{figure}[p]\centering
\includegraphics[width=6.4in]{figures/faraday_repeatability.pdf}
\caption{Repeated sidereal-phase measurements. Two simulated repeat measurements of the same sky and instrument at sidereal phase 0. Top: Q residuals, with parent points and zoom curves. Bottom: their difference divided by the predicted independent-repeat radiometer RMS, sqrt(2) times the individual RMS; gray is $\pm$1. The sky, smooth polarization, speckle and leakage are held fixed. Noise draws are independent; each cycle uses one eligible phase-bin exposure at 50\% duty, exported in seconds. Correlated power sampling uncertainties and independent mock observations are recorded, without counting channels as independent skies.}
\end{figure}
\begin{figure}[p]\centering
\includegraphics[width=6.4in]{figures/faraday_angular_and_fringe.pdf}
\caption{Angular and observed fringe spectra. Conditional unrotated, foreground and uniform-slab polarized angular spectra (top), and their Q fringe power through a common instrument and observing operator (bottom). All use the local 0.02 rad/m$^2$ dispersion illustration and 1 K intrinsic RMS at 30 MHz. The operator averages 256 sidereal bins, selects complete actual-site nightside bins and applies their normalized Fourier transform. The same mask and finite-bin average, with common harmonic cutoff ell=48, act on polarized covariance, deterministic Haslam I leakage and explicit loaded receiver/radiometer noise. Noise here uses a 25 kHz top-hat at 50\% duty per retained bin. The window mixes modes: low-ell input shape alone does not establish separability, and no angular-map sensitivity is inferred. Full complex fringe covariance and relation are exported. Angular power below the displayed twelve-decade range is omitted.}
\end{figure}


\clearpage
\section{Decision table and interpretation}

\begin{table}[htbp]
\centering\small
\caption{Full signed-column range: largest residual Q/U speckle RMS (K), speckle/noise and total-polarization/noise RMS ratios, and retained speckle power at phase 0 over the longest illustrative exposure. Maxima can come from different conditional cases. The full arrays retain each coordinate and unavailable exposure.}
\label{tab:decisions}
\begin{tabular}{llllllll}
\toprule
MHz & Type & Profile & Q (K) & U (K) & S/noise & Total/noise & Power kept\\
\midrule
30 & parent & f=0 & $0$ & $0$ & $0$ & $7.6\times10^{-7}$ & $\text{undefined}$\\
30 & parent & f=0.6 & $1.1\times10^{-4}$ & $9.8\times10^{-5}$ & $3.5\times10^{-4}$ & $2.9\times10^{-3}$ & $1.7\times10^{-5}$\\
30 & parent & f=1 & $1.1\times10^{-4}$ & $1.0\times10^{-4}$ & $3.5\times10^{-4}$ & $7.5\times10^{-3}$ & $1.3\times10^{-4}$\\
30 & parent & slab & $8.9\times10^{-7}$ & $7.8\times10^{-7}$ & $2.7\times10^{-6}$ & $5.9\times10^{-5}$ & $1.3\times10^{-4}$\\
30 & parent & Gaussian & $1.0\times10^{-7}$ & $9.0\times10^{-8}$ & $3.2\times10^{-7}$ & $2.6\times10^{-6}$ & $1.3\times10^{-4}$\\
30 & zoom & f=0 & $0$ & $0$ & $0$ & $8.8\times10^{-8}$ & $\text{undefined}$\\
30 & zoom & f=0.6 & $1.4\times10^{-4}$ & $1.2\times10^{-4}$ & $4.1\times10^{-5}$ & $3.4\times10^{-4}$ & $2.7\times10^{-5}$\\
30 & zoom & f=1 & $1.4\times10^{-4}$ & $1.3\times10^{-4}$ & $4.1\times10^{-5}$ & $8.7\times10^{-4}$ & $2.1\times10^{-4}$\\
30 & zoom & slab & $1.1\times10^{-6}$ & $9.7\times10^{-7}$ & $3.2\times10^{-7}$ & $6.8\times10^{-6}$ & $2.1\times10^{-4}$\\
30 & zoom & Gaussian & $1.3\times10^{-7}$ & $1.1\times10^{-7}$ & $3.7\times10^{-8}$ & $3.0\times10^{-7}$ & $2.1\times10^{-4}$\\
50 & parent & f=0 & $0$ & $0$ & $0$ & $9.6\times10^{-8}$ & $\text{undefined}$\\
50 & parent & f=0.6 & $2.3\times10^{-6}$ & $1.0\times10^{-6}$ & $4.7\times10^{-6}$ & $3.5\times10^{-5}$ & $4.0\times10^{-8}$\\
50 & parent & f=1 & $9.6\times10^{-6}$ & $4.1\times10^{-6}$ & $2.0\times10^{-5}$ & $9.9\times10^{-5}$ & $2.9\times10^{-7}$\\
50 & parent & slab & $2.2\times10^{-7}$ & $8.0\times10^{-8}$ & $4.4\times10^{-7}$ & $2.2\times10^{-6}$ & $2.9\times10^{-7}$\\
50 & parent & Gaussian & $2.4\times10^{-8}$ & $1.1\times10^{-8}$ & $4.9\times10^{-8}$ & $3.5\times10^{-7}$ & $3.1\times10^{-7}$\\
50 & zoom & f=0 & $0$ & $0$ & $0$ & $1.1\times10^{-8}$ & $\text{undefined}$\\
50 & zoom & f=0.6 & $2.9\times10^{-6}$ & $1.3\times10^{-6}$ & $5.4\times10^{-7}$ & $4.0\times10^{-6}$ & $6.2\times10^{-8}$\\
50 & zoom & f=1 & $1.2\times10^{-5}$ & $5.1\times10^{-6}$ & $2.3\times10^{-6}$ & $1.1\times10^{-5}$ & $4.6\times10^{-7}$\\
50 & zoom & slab & $2.7\times10^{-7}$ & $1.0\times10^{-7}$ & $5.1\times10^{-8}$ & $2.6\times10^{-7}$ & $4.6\times10^{-7}$\\
50 & zoom & Gaussian & $3.0\times10^{-8}$ & $1.4\times10^{-8}$ & $5.7\times10^{-9}$ & $4.1\times10^{-8}$ & $4.9\times10^{-7}$\\
\bottomrule
\end{tabular}
\end{table}

\begin{table}[htbp]
\centering\small
\caption{Phase-0 chromatic I-only residual RMS (K) through the common fit: as-built linear model, linear/PCHIP difference and worst ablation mismatch using the same as-built calibration. These model spreads are not measured beam-error probabilities.}
\label{tab:leakage}
\begin{tabular}{llllll}
\toprule
MHz & Type & Stokes & As-built & Interpolation & Mismatch\\
\midrule
30 & parent & Q & $2.6$ & $2.3$ & $5$\\
30 & parent & U & $4.5$ & $4.1$ & $7.2$\\
30 & zoom & Q & $2.9$ & $2.6$ & $5.5$\\
30 & zoom & U & $5$ & $4.5$ & $7.9$\\
50 & parent & Q & $1.1$ & $0.91$ & $0.91$\\
50 & parent & U & $0.31$ & $0.26$ & $0.26$\\
50 & zoom & Q & $1.2$ & $0.98$ & $0.98$\\
50 & zoom & U & $0.35$ & $0.28$ & $0.28$\\
\bottomrule
\end{tabular}
\end{table}

\begin{table}[htbp]
\centering\small
\caption{Conditional HII total-emissivity endpoints with assumed polarized fractions, converted to the reference frequency using the declared audit spectral index. Ratios rescale the same angular/depth model and its full emitting volume. These are not local measurements or a prior distribution.}
\label{tab:emissivity}
\begin{tabular}{llll}
\toprule
I at 74 MHz (K/pc) & Polarized fraction & P at 30 MHz (K/pc) & Max. total/noise\\
\midrule
$0.53$ & $0$ & $0$ & $0$\\
$0.53$ & $0.1$ & $0.61$ & $4.5\times10^{-3}$\\
$0.53$ & $0.7$ & $4.2$ & $0.032$\\
$1.5$ & $0$ & $0$ & $0$\\
$1.5$ & $0.1$ & $1.7$ & $0.013$\\
$1.5$ & $0.7$ & $12$ & $0.091$\\
$2.5$ & $0$ & $0$ & $0$\\
$2.5$ & $0.1$ & $2.9$ & $0.022$\\
$2.5$ & $0.7$ & $20$ & $0.15$\\
\bottomrule
\end{tabular}
\end{table}


The HII-based emissivity brackets multiply the full emitting normalization;
their largest total-polarization to noise RMS ratio is
\MaxEmissivityTotalNoise. The assumed polarized fractions, spectral
conversion and uniform local extension prevent interpreting that value as
a measured local upper bound. The polarization remains small enough relative
to I for the stated I-only noise approximation in these brackets.

The table reports Q and U separately and distinguishes power transmission
from amplitude ratios. Beam ablations and interpolation choices quantify
model differences, not calibrated beam-error probabilities. A correlated
mean-square statistic uses its full covariance sampling variance;
channels are not counted as independent sky realizations. Independent
simulated observations pass through the calibrated-spectrum input path,
and paired I-only controls share the same noise draw.

Angular spectra are compared with leakage only after the same instrument,
finite time bins, actual-site night mask and Fourier operator. Masked
fringe power is an observable comparison; it is not a reconstructed angular
map or a map-level sensitivity claim. An approximately white input angular
spectrum alone does not establish separation.

\section{Approximations and next steps for mission data}

The delivered predictions are conditional on the local-medium and receiver
inputs, the intrinsic angular spectrum and low-frequency turnover, fixed
narrow-band signal kernels, interpolated chromatic leakage, and independent
radiometer averages. A measured spectral residual or a repeated sky pattern
does not by itself identify Faraday rotation. The observing operator,
beam/gain uncertainty and competing sky terms must accompany a physical
model comparison.

When LuSEE-Night data become available, supply their calibrated Stokes
spectra with the actual response, flags, phase bins and integration, then
regenerate the common estimator and noise prediction. Compare residual
variance, named-bin frequency correlations and same-phase repeat differences
with conditional predictive distributions. A discrepancy or an upper bound
is a valid outcome. Converting a floor into a distance or field inference
requires independent information that breaks the emissivity/depth/coherence
degeneracy.

\section{Provenance and regeneration}

The report source and products generator were committed before regeneration. The build used a clean detached worktree at source revision \texttt{56975a30e7b5}. Numerical records contain complete source and dependency manifests; the report manifest checksums every record, array and figure it reads. Large files remain regenerable. Existing unrelated notebook edits in the shared workspace are preserved and are not report inputs.


The external paper has not been modified. The earlier instrument and
point-source material remains on the \texttt{audit-2026-08-18} tag and in the
existing analysis scripts. The superseded detection forecast is not a
required result of this calculation. Regenerate with
\nolinkurl{scripts/run_paper_implementation.sh}; use its help for stages
and the data-free smoke mode. All large numerical arrays are checkpointed
under \nolinkurl{generated_data/cache/}; figure and report manifests retain
their complete hashes.

\end{document}
